\newcommand{\CaseID}{bobcat\_2020}
\newcommand{\CaseTitle}{2020 Bobcat fire reconstruction against GOFER observations}
\newcommand{\FireName}{Bobcat}
\newcommand{\FireYear}{2020}
\newcommand{\EventStart}{6 September 2020 at 19:00~UTC}
\newcommand{\EventEnd}{24 September 2020 at 20:00~UTC}
\newcommand{\DurationHours}{433}
\newcommand{\CRSName}{UTM zone 11N (EPSG:32611)}
\newcommand{\FuelVersionText}{LANDFIRE release 2.1.0}
\newcommand{\FireSummary}{The supplied GOFER event record contains 433 hourly perimeters, a final area of 583.864~km$^2$ and an official size of 115,997~ac. The configuration uses 3 timed ignition point(s).}
\newcommand{\ConfigurationText}{The supplied simulation runs for 433~h, with an initial timestep of 1~s, target Courant--Friedrichs--Lewy number of 0.2, 434 hourly weather fields and 3 ignition point(s). Arrival time, spread rate, fireline intensity and fire-size statistics are requested at the supplied dump interval. Crown-fire model 1 and firebrand spotting are enabled. Unspecified settings retain the inspected source defaults, including fixed random seed 2024, foliar moisture 90\% and maximum timestep 600~s. A single simulation member is configured.}
\newcommand{\WeatherOffsetText}{Simulation and weather time zero is 0~h after the GOES ignition hour; comparisons apply this offset explicitly.}
\newcommand{\ElmfireRevision}{8981f14e0170}
\newcommand{\ElmfireBranch}{yiren}
